\section{Conclusion}
\label{sec:conclusion}

We asked whether AI agents can communicate covertly without ever sharing a secret key. They can: an interactive, pseudorandom key exchange hidden in sampling randomness works, and we demonstrated it through a real LM with 10/10 key agreements and no detectable signal for an equally informed auditor. What agents cannot do without is (a) prior common knowledge of the protocol, (b) interaction with private coins, (c) entropy in honest messages, and (d) a passive monitor. Tacit ``secret languages'' give zero secrecy against a monitor with the same model and knowledge, both in theory and in our experiments (48\% vs.\ 50\% recovery). We recommend that overseers run monitors with the same weights and context as the agents, and prefer active defenses such as light resampling, impersonation probes, and monitor-supplied sampling randomness over purely passive auditing.

Several directions remain open. First, repeating \eone with frontier models that are allowed to read the keyless-protocol papers would test whether publication turns a protocol into a focal point that agents run unprompted. Second, synchronization-robust gated channels (marker codes, per-message re-anchoring) may recover robustness against resampling wardens, and measuring how much would sharpen the defense. Third, \propref{prop:learn} deserves a formal proof under a precise undetectability notion. Finally, meaning-level channels should be measured under paraphrase.
